\section{相关工作}
\label{sec:related}

\paragraph{查询级路由与级联。}
早期研究将模型选择视为对完整查询在质量和成本之间进行的推理时决策。FrugalGPT 学习异构 API 的级联策略以降低成本~\citep{chen2024frugalgpt}。AutoMix 先查询较小模型，通过自验证估计可靠性，并在需要时升级~\citep{aggarwal2023automix}。Hybrid LLM 根据预测难度在本地模型与云端模型间路由~\citep{ding2024hybrid}。RouteLLM 从偏好数据训练路由器，并研究跨模型对的迁移~\citep{ong2024routellm}。一篇近期综述将共同目标形式化为在成本预算约束下选择模型~\citep{varangotreille2025doing}。这些方法有效地抽象了完整查询上的模型选择，但即便工具输出和局部状态已塑造前缀，多轮智能体轨迹中的中间调用仍缺乏充分研究。

\paragraph{路由器基准。}
RouterBench 收集了跨多样化单提示词任务的逾 40.5 万条推理结果，提供了首个用于比较路由算法的大规模测试平台~\citep{hu2024routerbench}。LLMRouterBench 将规模扩展至逾 40 万个实例、21 个数据集和 33 个模型，并采用统一的性能与成本指标~\citep{li2026llmrouterbench}；尽管其中包含 500 个 SWE-Bench issue，SWE-Bench 被配置为无智能体脚手架、无中间可路由调用的单查询 $0$-shot \textsc{Pass@1} 任务。RouterArena 增加了带难度分级和自动排行榜的开放比较平台~\citep{lu2025routerarena}。Triage 使用代码健康度特征，对每个 SWE-bench Lite issue 作一次层级决策~\citep{madeyski2026triage}。这些基准对于比较查询级路由器很有价值；但即使纳入 SWE-Bench 等多步任务，路由仍被简化为从预先计算的结果矩阵中为每个 issue 选择一个模型，无法逐一将 issue 内部的中间检索、分析和调试调用路由到每一步足够便宜的模型。

\paragraph{步级路由。}
TRIM 通过识别关键推理步骤并将其升级到更大模型，研究多步数学推理的逐步路由；它表明有选择地逐步升级可以在降低平均成本的同时保持准确率~\citep{kapoor2026trim}。但数学推理是相对自包含的场景，缺少真实智能体部署的核心要素：工具输出、检索上下文、shell 轨迹和代码状态。TRIM 还从假设的推理轨迹导出标签，而未端到端执行降级后的轨迹来验证任务成功。对于可执行的长上下文智能体工作负载，核心标注问题依然存在：在给定当前前缀时，哪个模型层级既足够便宜，又仍能保持最终任务解决？
